%============================================================
% BAB 2 — KAJIAN PUSTAKA
%============================================================
\chapter{KAJIAN PUSTAKA}

Awali pembahasan pada bab ini dengan mengutarakan \textit{Problem} pada
\textit{Problem Domain} dari proyek akhir. Setelah itu, penjelasan tersebut diiringi
dengan teori-teori penunjang pada bidang pengetahuan yang terkait dengan
\textit{Problem Domain} dan \textit{Problem}. Untuk lebih menguatkan klaim orisinalitas
proyek akhir, perlu diberikan ulasan tentang penelitian-penelitian terkait yang juga
bertujuan menyelesaikan \textit{Problem} tersebut.

\section{DESKRIPSI PERMASALAHAN}

Deskripsikan dengan jelas dan detail dari permasalahan yang ingin diselesaikan pada
proyek akhir. Permasalahan berisi penjelasan dari \textit{Problem} yang termuat pada
judul kegiatan. Deskripsi masalah sebaiknya dituliskan dengan gaya bahasa deskriptif.
Deskripsi masalah boleh memuat gambar, tabel, dan skema tertentu untuk mengilustrasikan
permasalahan.

\section{TEORI PENUNJANG}

Uraikan dengan jelas dasar teori yang menunjang penelitian proyek akhir. Teori penunjang
menguraikan dasar-dasar teori, temuan, dan bahan dari pustaka ilmiah lain yang dijadikan
landasan untuk melakukan proyek akhir yang diusulkan. Uraian dalam Teori Penunjang
menjadi landasan untuk menyusun kerangka atau konsep yang akan digunakan dalam proyek
akhir. Jangan lupa untuk menyebutkan semua kutipan dengan rujukan yang jelas seperti
ini \cite{buku_contoh}.

\subsection{Sub-Teori Pertama}

Uraikan sub-teori pertama di sini.

\subsection{Sub-Teori Kedua}

Uraikan sub-teori kedua di sini.

\section{PENELITIAN TERKAIT}

Penelitian terkait memuat hasil penelitian pihak lain yang mempunyai \textit{Problem}
yang sama dengan penelitian kita, tetapi dengan menggunakan \textit{Uniqueness} yang
berbeda. Ceritakan bagaimana penelitian-penelitian terkait telah mencoba menyelesaikan
permasalahan yang sama, dengan cara mereka masing-masing, yang ditunjukkan dengan
kutipan terhadap pustaka. Penelitian terkait yang baik melibatkan kajian pustaka yang
relevan dan terpercaya dari jurnal ilmiah internasional ataupun nasional, presentasi
ilmiah internasional ataupun nasional, dan buku atau catatan rujukan ilmiah. Penulis
harus mencantumkan sumber referensi pada daftar pustaka manakala melakukan rujukan dan
kutipan dari pihak lain secara jujur dan benar seperti ini \cite{jurnal_contoh}.
Pencantuman sumber referensi perlu dilakukan baik terhadap kutipan langsung ataupun
kutipan tidak langsung (parafrase).

Untuk kutipan langsung dan pendek (1--2 baris), cara penulisan rujukan dilakukan dengan
memberikan tanda petik ganda di awal dan akhir kutipan, ditulis miring, kemudian diiringi
sumber referensi, seperti ini \textit{``Ini contoh penulisan rujukan untuk kutipan
langsung dan pendek''}\cite{buku_contoh,jurnal_contoh}. Sedangkan untuk kutipan langsung
dan panjang (lebih dari 2 baris), penulis dapat menuliskannya seperti di bawah ini.

\begin{kutipanpanjang}
  ``Ini contoh penulisan rujukan untuk kutipan langsung dan panjang, ditulis pada
  paragraf terpisah dengan memberikan tanda petik ganda di awal dan akhir kutipan,
  ukuran font 10 point dan margin kanan kiri yang masuk 10 mm dari batas penulisan,
  kemudian diiringi dengan sumber referensi pada daftar pustaka.'' \cite{buku_contoh}
\end{kutipanpanjang}

Untuk kutipan tidak langsung (parafrase), penulis dapat menuliskan sumber referensi
setelah kalimat parafrase selesai, seperti ini \cite{online_contoh}.